\chapter{Using what we already know: priors, transfer, and the price of a trial}
\label{sec:transfer}

A tuner that starts every study from ignorance wastes its first dozen trials
rediscovering folklore: learning rates live near $3 \times 10^{-4}$, discount
factors near $0.99$, dropout rarely helps reinforcement learning. Every GP in
Chapter~\ref{sec:bayesian} began with a zero-mean prior that believes
nothing; every project lead we serve believes a great deal, and is usually
right. We tune the same families of models over and over, so encoded
experience is the cheapest compute we own. The literature offers three
grades of it, in increasing order of machinery, plus a fourth option that
replaces search entirely at large scale.

\section{Expert priors over the optimum}

The lightest grade lets a human say where the optimum probably lies. In
$\pi$BO, the user supplies a prior density $\pi(\config)$ over the
\emph{location of the best configuration} (not over function values, which
is the GP's job), and the acquisition function is multiplied by that
density raised to a power that decays with the number of observations,
\begin{equation}
\alpha_n^{\pi}(\config)
\;=\; \alpha_n(\config) \cdot \pi(\config)^{\beta / n},
\label{eq:pibo}
\end{equation}
where $\alpha_n$ is any standard acquisition (EI in the paper), $n$ counts
observations, and $\beta$ sets how long the advice matters
\citep{hvarfner2022pibo}.

The design reads directly off the formula. Early
on ($n$ small) the exponent is large, the prior term dominates, and the
search concentrates where the expert pointed; as evidence accumulates the
exponent shrinks toward zero, $\pi(\config)^{\beta/n} \to 1$ for every
$\config$ in the prior's support, and the acquisition converges back to
its data-driven self on that support (Figure~\ref{fig:pibo}). Points
assigned zero prior probability stay at zero. The decay is what makes
the mechanism safe: the authors prove it preserves standard convergence
behavior, and show large early gains when the prior is good and mild
losses when it is wrong, since a wrong prior is overruled at the rate
evidence arrives. A confidently wrong expert costs the study its opening
moves, not its conclusion.

\begin{figure}[t]
\centering
\begin{tikzpicture}[scale=1.0, font=\small]
\draw[-{Latex}] (0,0) -- (9.3,0);
\node[anchor=north] at (7.9,-0.12) {learning rate (log scale)};
% prior density
\draw[thick, densely dashed]
  (0.5,0.1) .. controls (2.8,0.12) and (3.5,1.8) .. (4.1,1.85)
  .. controls (4.7,1.8) and (5.4,0.15) .. (8.9,0.07);
% early acquisition: pulled toward prior
\draw[thick, blue!60!black]
  (0.5,0.25) .. controls (2.6,0.4) and (3.4,2.3) .. (4.2,2.4)
  .. controls (5.0,2.3) and (5.6,0.7) .. (8.9,0.3);
% late acquisition: data-driven peak elsewhere
\draw[thick, red!60!black]
  (0.5,0.15) .. controls (4.0,0.25) and (6.0,0.4) .. (6.8,1.5)
  .. controls (7.3,2.05) and (7.7,1.35) .. (8.9,0.25);
% legend, top left, clear of all three curves
\node[blue!60!black, anchor=west, font=\footnotesize] at (0.35,3.25)
  {$\alpha^{\pi}_{n}$, $n$ small: the prior steers};
\node[red!60!black, anchor=west, font=\footnotesize] at (0.35,2.9)
  {$\alpha^{\pi}_{n}$, $n$ large: exponent $\beta/n$ decayed, data decides};
\node[anchor=west, font=\footnotesize] at (0.35,2.55)
  {expert prior $\pi$ (dashed)};
\end{tikzpicture}
\caption{The $\pi$BO mechanism~\eqref{eq:pibo} over a study's lifetime
(schematic). The expert's prior (dashed) multiplies the acquisition with
an exponent that decays as observations arrive: early proposals
concentrate where the expert pointed, late proposals go where the data
says, and a wrong prior is forgotten rather than obeyed.}
\label{fig:pibo}
\end{figure}

PriorBand carries the same idea into the multi-fidelity world of
Chapter~\ref{sec:multifidelity}, where the object being steered is not
one acquisition but a stream of configurations entering ASHA-style rungs.
Each new configuration is drawn from a portfolio mixing three samplers:
uniform exploration, the user prior, and local perturbations of the
current incumbent, with the mixing weights depending on the rung. The
prior is trusted more at higher rungs, where evaluations are expensive
and the expert's opinion matters most, and the incumbent component is
switched on only once enough evidence has accumulated, which is what
protects the method against a confidently wrong prior
\citep{mallik2023priorband}. Priors of this kind should be a first-class
field in our search-space schema from day one: every project lead we
serve has them, and both mechanisms are simple to implement over
components we are building anyway. The prior is also the natural home
for the folklore this chapter opened with; writing
$\pi(\config)$ down once beats re-teaching it to every study.

\section{Transfer from our own past studies}

The second grade learns from previous tuning studies instead of asking a
human. The robust end of this spectrum is unglamorous and effective:
initialize a new study with configurations that ranked well on similar
past tasks \citep{feurer2015metalearning}, or combine surrogates fit to
past tasks into an ensemble whose weights reflect how well each past
task's surrogate predicts the current one, comparing on ranks so that
tasks with different loss scales can vote together
\citep{feurer2018rgpe}. Because our product will keep every trial of
every study in one store (Chapter~\ref{sec:architecture}), the warm
start costs a database query: the run store is Figure~\ref{fig:loop}'s
bottom box earning its place.

The research frontier goes further and meta-learns the surrogate itself,
and the vehicle deserves a proper introduction because it has appeared
twice already (ifBO in Chapter~\ref{sec:multifidelity}) and will appear
again. A \emph{prior-fitted network} (PFN) is a transformer trained once,
offline, on millions of synthetic datasets drawn from a chosen prior over
functions. At use time it is shown the current study's observed
(configuration, score) pairs as context, and it outputs a predictive
distribution for new configurations in a single forward pass.

Why that works is the part worth pausing on. Training on draws from
the prior teaches the network to approximate the posterior predictive
distribution under that prior, so the expensive part of Bayesian
inference is paid once, by the pre-trainer, rather than at every step
of every study. Where the GP of Section~\ref{sec:gp} refits
\eqref{eq:mll} inside the loop, a PFN just reads its context. The
model class changed; the economics changed with it.

PFNs4BO demonstrates the swap directly, pre-training PFNs that mimic
GP and Bayesian-neural-network surrogates and slotting them into
standard BO \citep{muller2023pfns4bo}.

OptFormer showed the scale endgame: a transformer trained on hundreds
of thousands of tuning trajectories from Google's Vizier service
imitates several tuners and improves on them when combined with
expected improvement \citep{chen2022optformer}.

The same idea specialized to learning curves gives ifBO's freeze-thaw
surrogate (Chapter~\ref{sec:multifidelity}), and the line is still
moving, with in-context amortization now reaching the acquisition
function itself \citep{rakotoarison2026alphapfn}.

These pre-trained surrogates are the most credible path to tuners
that get smarter as our run store grows, and one of them (ifBO) ships
with usable code today. We rate them a second-wave investment:
valuable, but only after the plain warm start is in place and
earning.

\section{Cost-aware search}

The third grade treats the price of a trial as part of the decision.
Trials rarely cost the same: in the running example, doubling network
width or unroll length can quadruple wall-clock, so two candidates with
equal expected improvement can differ fourfold in what they cost to
evaluate, and an acquisition that ignores this will happily spend its
budget in the expensive corner of the space. The classical repair prices
the acquisition in the right currency: improvement per unit cost,
$\alpha_{\mathrm{EI}}(\config) / c(\config)$, with $c$ a model of trial
cost, already present in \citet{snoek2012practical} as expected
improvement per second. CArBO refines it with an explicit cheap-first
initialization and a cost-cooling schedule, dividing by
$c(\config)^{\alpha}$ with an exponent that anneals from one toward
zero over the budget, so that early exploration favors cheap points and
the endgame, where only the best candidates remain, ignores cost
\citep{lee2020carbo}; the annealing acknowledges that a bargain matters
when we are learning the space and stops mattering when we are
choosing the winner.

CARBS pushes cost-awareness to its logical end for scaling studies. It
maintains a Pareto front in the plane of cost versus performance, the
machinery of Chapter~\ref{sec:moo} with cost promoted to a full
objective, and searches locally around that front: separate Gaussian
processes model performance, cost, and the front's shape, candidates are
Gaussian perturbations of current Pareto members, and the acquisition is
expected improvement relative to the front weighted by expected cost, so
the method tunes hyperparameters \emph{while} the model grows and emits
a scaling law for every hyperparameter as a byproduct. Its authors
demonstrate it by reproducing the Chinchilla scaling result for language
models and by solving the full ProcGen benchmark with a tuned but
otherwise plain PPO \citep{fetterman2023carbs}. For our fixed-size
studies, EI-per-cost with cost cooling is the sensible default; a
CARBS-style cost frontier becomes interesting the day we tune models
meant to be scaled up afterward.

\section{When not to search at all: transfer across scale}

For large networks there is a fourth option that is not search. The
maximal update parametrization ($\mu$P) rescales initialization and
per-layer learning rates so that the size of each layer's update stays
controlled as network width grows; under that parametrization the optimal
values of many hyperparameters become approximately invariant to width,
so one tunes a small proxy model and transfers the result to the wide
target ($\mu$Transfer). Its authors validated the recipe by
outperforming the published BERT-large and GPT-3 6.7B settings while
spending a few percent of one pretraining run on tuning
\citep{yang2022mutransfer}. The claim's boundary is part of the result:
the theory covers width; transfer across depth, batch size, and training
horizon is empirical, with stated failure cases (regularization
hyperparameters do not transfer, and depth transfer breaks for
post-layer-norm transformers), and follow-up parametrizations have since
extended the depth story \citep{dey2025completep}.

More recent LLM-era work makes the third regime of
Figure~\ref{fig:regimes} explicit, fitting power laws for optimal
hyperparameters as functions of scale and turning the largest tuning
problems into curve lookups: compute-dependent learning-rate and
batch-size laws appear in the DeepSeek LLM report \citep{deepseek2024llm},
learning-rate dependence on the token horizon has its own fitted law
\citep{bjorck2025tokenhorizon}, weight decay and batch size get the same
treatment \citep{bergsma2025powerlines}, and a dedicated preprint fits
joint optimum surfaces across model and data size \citep{li2025steplaw}.

We track this line in Chapter~\ref{sec:roadmap} as guidance to encode
rather than an algorithm to implement: the laws enter our product as
priors and sanity checks, not as a searcher.

\begin{figure}[t]
\centering
\begin{tikzpicture}[font=\small]
\draw[-{Latex}] (0,0) -- (12.6,0) node[above left=0.6mm and 0mm, font=\small\itshape] {model scale};
\foreach \x in {4.2,8.4} \draw[dotted] (\x,0) -- (\x,2.6);
\node[align=center] at (2.1,1.7) {\textbf{search directly}\\Chapters~\ref{sec:bayesian}--\ref{sec:moo}:\\trials are affordable,\\evidence beats folklore};
\node[align=center] at (6.3,1.7) {\textbf{search a proxy, transfer}\\$\mu$Transfer: tune small,\\move to the wide model\\\citep{yang2022mutransfer}};
\node[align=center] at (10.5,1.7) {\textbf{look up fitted laws}\\scaling rules for LR, batch,\\weight decay; one model,\\zero search trials};
\node[anchor=north, font=\footnotesize] at (2.1,-0.35) {our RL, samplers, HNNs};
\node[anchor=north, font=\footnotesize] at (6.3,-0.35) {medium pretraining};
\node[anchor=north, font=\footnotesize] at (10.5,-0.35) {frontier-scale pretraining};
\end{tikzpicture}
\caption{The three tuning regimes by model scale. Everything before this
chapter lives in the left band. As trials become unaffordable, search
moves onto small proxies whose optima provably or empirically transfer,
and at the largest scales it gives way to fitted scaling laws. A good
tuner should know which band a study is in and say so.}
\label{fig:regimes}
\end{figure}

The strategic point for the product is the figure's caption: at small
scale we search, at medium scale we search on proxies and transfer, and
at large scale search gives way to parametrization and extrapolation. A
good tuner should know which regime it is in and say so.

\section{The evidence, and how the failure mode was tested}
\label{sec:prior-evidence}

Every claim in this chapter has the same shape: knowledge from outside
the study, an expert's hunch, a past study, a smaller model, someone
else's fitted law, can safely be spent inside it. For claims of that
shape the success story is cheap to produce and nearly beside the
point, because of course good advice helps. The evidence that matters
is what happens when the imported knowledge is wrong, and the reason
this chapter's verdicts lean confident is that its sources, unusually,
tested exactly that.

The $\pi$BO experiments are built around controllable advice quality.
On a toy function and two surrogate tuning tasks, where the true
optimum is known and priors of any desired quality can be
manufactured, the authors ran the spectrum from accurate to
``purposefully detrimental,'' the wrong prior being a narrow density
parked in the worst region of the space \citep{hvarfner2022pibo}.

The findings match the mechanism read off
formula~\eqref{eq:pibo}: strong priors buy large early gains, and
under sabotage the method climbs back to roughly the regret of its
prior-free baseline, while the baseline that trusts the prior only at
initialization keeps paying for it. The trade-off is stated rather
than hidden: a larger $\beta$ extracts more from good advice and
forgives bad advice more slowly.

The upside case is demonstrated
separately on a deep-learning pipeline, a 12.5-fold time-to-accuracy
speedup over vanilla BO and a new best validation accuracy for that
pipeline. Author-run experiments, as usual in this document's
sources; what tempers the caution here is that the safety half of the
claim is also a theorem (the exponent $\beta/n$ dies as evidence
arrives, dragging the acquisition back to its data-driven self), so
the experiment confirms an argument the reader can check rather than
substituting for one.

PriorBand's evaluation had to build its own arena first, a benchmark
suite in which priors are part of the task definition rather than an
afterthought: synthetic multi-fidelity surfaces plus image and
language tasks derived from recorded large-model training runs, each
equipped with good and deliberately bad priors, run at the small
budgets that define deep-learning tuning
\citep{mallik2023priorband}. The motivating ablation is the honest
part: naively biasing Hyperband's sampler toward the prior wins when
the advice is good and burns exactly proportional budget when it is
bad, while a half-random hedge is robust but never competitive, and
the rung-dependent trust schedule is the repair that keeps both
properties. One design detail is worth copying into our own suite:
every prior-using method in the comparison had the prior's mode
evaluated first, so nobody wins merely by being handed the expert's
best guess. Across tasks and prior qualities, PriorBand ranked most
consistently at the top.

The transfer-across-scale evidence is differently shaped, and
stronger for it. $\mu$Transfer's headline test pitted the recipe
against external baselines its authors did not set: the published,
expert-tuned settings of BERT-large and GPT-3 6.7B, beaten while
spending a few percent of one pretraining run on proxy-model tuning
\citep{yang2022mutransfer}, with the non-transfers (regularization
strengths, depth under post-layer-norm) reported in the same paper.
CARBS makes its case the same way, by reproducing a famous external
result, the Chinchilla compute-optimal scaling law, and by solving
the full ProcGen benchmark with an otherwise plain PPO
\citep{fetterman2023carbs}. Replication of results the authors do
not own is the closest thing this literature has to the
adversarial arena of Section~\ref{sec:bo-evidence}, and both
transfer claims clear it.

The last tier is the one we deliberately did not audit. The
scaling-law papers of the previous section report fitted power laws
whose estimation we have not checked, and nothing in our product is
allowed to depend on their exactness: they enter as priors and
sanity checks, which is weight their evidentiary status can carry
even if a particular fit is off. The demotion is a statement about
our audit, not an accusation about their work.

\section{The ordering, argued from four angles}
\label{sec:prior-angles}

The chapter's verdict is an ordering: priors and warm starts first,
cost-aware acquisition alongside, pre-trained surrogates adopted
rather than trained, scaling knowledge encoded rather than
implemented. Four angles converge on it.

From mechanism: the safety properties that make priors acceptable as
a default are structural, not empirical. The decay exponent in
\eqref{eq:pibo} guarantees the expert's influence is transient;
PriorBand's incumbent component stays off until evidence justifies
it; in both, the worst case of bad advice is a study's opening
moves. Features whose failure modes are bounded by construction can
be on by default; features whose failure modes are empirical
questions cannot. That one sentence sorts most of the decision box.

From evidence: the sources tested the sabotage case, not only the
success case, and the transfer methods validated against external
baselines their authors did not control. Where our own audit is
thinner, the scaling-law fits, the verdict steps down to
encode-as-guidance in lockstep. Nothing here asks the reader to
extend more trust than the checked evidence covers.

From economics: encoded experience is the cheapest compute we own.
A written-down prior converts folklore into an asset that every
study inherits for free; the run store turns past studies' sunk
compute into a warm start that costs a query; and EI-per-cost simply
prices candidates in the currency the cluster actually charges,
improvement per unit cost rather than improvement per trial. The
whole chapter is arbitrage on knowledge the organization has already
paid for once.

From operations: the entire first wave is days of work riding
components that earlier chapters build anyway, a multiplier on the
acquisition, a query over the run store, a divisor on the
acquisition, and the single costly commitment, making the prior a
first-class field of the search-space schema, costs nearly nothing
on day one and a migration if retrofitted later. Adopting ifBO's
published surrogate rather than training our own keeps a research
program off the product's critical path.

And the falsification criterion, stated in advance: the validation
suite reruns the pinned tasks under three conditions, no prior, the
folklore prior, and a deliberately wrong prior in $\pi$BO's sense,
at our routine budgets. The prior-weighted searcher must beat its
prior-free self under the folklore prior and must recover to the
prior-free final, within noise, under the wrong one; if bad advice
costs more than the opening moves on our workloads, prior weighting
is demoted from default to opt-in. The warm start faces its own
test the day the run store holds enough studies: measured trial
savings on a new study of a familiar family, or it stays a query we
offer rather than a step we perform.

\begin{decisions}
\noindent\textbf{Decisions from this chapter.} Ordered by payoff per
week of engineering, which is the honest ranking here.
\begin{itemize}[leftmargin=1.3em, itemsep=2pt]
\item \textbf{$\pi$BO's prior-weighted acquisition} (include; days).
Formula~\eqref{eq:pibo} is a few lines over the Tier 0 GP searcher,
every project lead has a prior to give it, and its built-in decay
means bad advice costs a study its opening moves rather than its
conclusion, \emph{provided the prior does not assign zero probability
to the true region}. The search-space schema should therefore require
a lower-bounded mixture or explicitly nonzero support everywhere, so
that wrong advice can be overruled by evidence \citep{hvarfner2022pibo}.
Highest value density in the document under that guard.
\item \textbf{PriorBand} (include; one to two weeks). The same idea
where our savings actually live, the multi-fidelity rungs; the
rung-dependent portfolio is what keeps a confidently wrong prior from
doing damage \citep{mallik2023priorband}.
\item \textbf{Warm starts from the run store} (include; days). A
ranked query over past studies \citep{feurer2015metalearning}. The
prerequisite is the run store itself, which Tier 0 builds regardless.
\item \textbf{RGPE-style surrogate ensembles} (second wave). Real
gains, but only once the store holds enough completed studies to
ensemble; building it before the data exists would be decoration.
\item \textbf{Pre-trained surrogates (PFN family)} (adopt one, build
none). We take ifBO's published pre-trained model as our
learning-curve surrogate (Chapter~\ref{sec:multifidelity}) and train
nothing ourselves; OptFormer-scale training runs are a research
program, not a package feature.
\item \textbf{EI-per-cost with cost cooling} (include; days). One
division and one annealing exponent on top of the GP searcher,
\emph{dividing by a predicted cost} rather than the realized cost the
executor measures \citep{lee2020carbo}. A second GP fits log
wall-clock time to configuration coordinates, refitting on the same
cadence as the performance surrogate, with failed and killed trials
entered as censored observations. The roadmap details the cold-start
and calibration-logging protocol.
\item \textbf{CARBS-style cost frontier} (hold). Becomes relevant the
day a project tunes a model that is destined to be scaled up; none is
on the books today \citep{fetterman2023carbs}.
\item \textbf{$\mu$P and the scaling-law literature} (encode, do not
implement). These enter the package as priors, defaults, and sanity
checks on the regime map of Figure~\ref{fig:regimes}, not as a
searcher; at our current scales they advise rather than replace
search.
\end{itemize}
\end{decisions}

\section{What would overturn these verdicts}

The $\pi$BO prior weight's default-on status depends on V11: if
gains under a good prior do not materialize, or if recovery from a
wrong prior takes longer than the no-prior baseline, the feature is
demoted to opt-in. PriorBand's inclusion holds as long as its ASHA
dependency does; if ASHA is demoted, PriorBand's rung-sampling
portfolio loses its substrate. Warm starts (V12) mature as the store
fills; if trial savings remain undetectable after the store holds
dozens of studies, the query surface is re-examined. EI-per-cost's
inclusion assumes executor cost metering is cheap; if metering
overhead becomes visible in profiling, the feature becomes opt-in.
RGPE's second-wave hold is lifted when an ensemble that cheap proves
itself on public suites, at which point the evidence section is
updated and a new decision box opened.
